% TOP_PROBLEM: {"id":30007658,"problem_number":"LOCAL-30007658","title":"Formalization without destroying livelihoods","category_id":21,"category":{"id":21,"name":"economics","display_name":"Economics","description":"Open economic theory statements, published research agendas, and proposed scoped specifications; question provenance and historical priority are qualified per record.","slug":"economics","order_index":21,"created_at":"2026-10-08T00:00:00Z"},"status":"research_question","external_url":null,"legacy_ids":[],"published":false,"statement":"In the explicitly specified setting: Which combinations of enforcement, registration incentives, public services, and social insurance make formalization raise productivity and welfare rather than merely displace informal work? The mathematical statement above fixes the target, assumptions and scope of this catalogue formulation.","economics":{"original_id":147,"field":"Development and poverty","kind":"design","formalization_basis":"adapted_research_question","source_status":"explicit_open","priority_status":"earliest_found","priority_note":"Earliest directly inspected qualifying passage in this audit; earlier origins were not established. A review or research-agenda statement does not imply original priority.","earliest_verified_reference":{"authors":["Gabriel Ulyssea","Matteo Bobba","Lucie Gadenne","Mariaflavia Harari"],"title":"Informality, VoxDevLit 6(2)","year":2025,"url":"https://voxdev.org/voxdevlit/informality-issue-2/final-remarks-informality","doi":null,"locator":"Chapter 6, Final Remarks, paragraphs 3–6","evidence_summary":"Firm dynamics and workers’ formal-informal transitions are explicitly unresolved, including stepping-stone versus trap mechanisms."},"current_reference":{"authors":["Gabriel Ulyssea","Matteo Bobba","Lucie Gadenne","Mariaflavia Harari"],"title":"Informality, VoxDevLit 6(2)","year":2025,"url":"https://voxdev.org/voxdevlit/informality-issue-2/final-remarks-informality","doi":null,"locator":"Chapter 6, Final Remarks, paragraphs 3–6","evidence_summary":"Firm dynamics and workers’ formal-informal transitions are explicitly unresolved, including stepping-stone versus trap mechanisms."},"formalization_note":"The cited passage establishes a research direction, not this exact mathematical statement. The notation, population, policy menu, assumptions, and estimands are an original adaptation for this catalogue; no universal unsolved theorem or source priority is claimed."},"statusQualification":"Published question or research agenda verified at the cited locator. The mathematical specification is adapted for this catalogue and includes additional assumptions; source publication does not establish first-ever priority or certify this exact model remains unsolved. The cited passage establishes a research direction, not this exact mathematical statement. The notation, population, policy menu, assumptions, and estimands are an original adaptation for this catalogue; no universal unsolved theorem or source priority is claimed.","statusReviewedAt":"2026-10-08"}
% FIRST_POSTED: 2026-10-08
% ENTRY_KIND: statement_only
% !TeX program = lualatex
\documentclass[11pt]{article}
\usepackage{fontspec}
\usepackage[a4paper,margin=25mm]{geometry}
\usepackage{amsmath,amssymb,mathtools}
\usepackage{xurl}
\usepackage[hidelinks]{hyperref}
\newcommand{\sref}[2]{\hyperlink{source:#1:#2}{[S#2]}}
\setlength{\emergencystretch}{3em}
\begin{document}
% problemId:30007658
\section*{Formalization without destroying livelihoods}\label{30007658}
\subsection{Definitions and mathematical statement}
Fix a labor and product market with informal firms and workers. Policy \(p\) specifies registration costs, tax and labor enforcement, access to services, and social-insurance eligibility. Let \(F_{it}(p)\) indicate formal status, \(Y_{it}(p)\) disposable earnings, and \(U_{it}(p)\) a declared welfare measure including work risk and benefits. Let \(E_t(p)\) denote firm entry and exit. Estimate dynamic effects \(E[F_{iT}(p)-F_{iT}(0)]\) and \(E[U_{iT}(p)-U_{iT}(0)]\), including workers and entrepreneurs initially informal and people displaced by enforcement. Determine whether informal activity serves as an entry route toward productive formal businesses and jobs or sustains persistent segmentation. Use linked longitudinal observations and policy variation to separate selection from causal transitions. Registration alone is not the welfare objective: a policy can increase formal shares by destroying marginal livelihoods. Counterfactuals must therefore incorporate exit, wages, prices, migration, and fiscal use of new revenue. Compare feasible policy bundles under explicit budget and enforcement constraints. The cited living review expressly identifies dynamic firm decisions and worker formal--informal transitions as research gaps; this welfare-centered intervention formulation extends those passages without asserting a universally beneficial formalization policy.

\par\textbf{Formulation and scope.} The cited passage establishes a research direction, not this exact mathematical statement. The notation, population, policy menu, assumptions, and estimands are an original adaptation for this catalogue; no universal unsolved theorem or source priority is claimed.

\par\textbf{Open-question provenance.} An explicit question or research agenda was verified at the source locator. This does not by itself certify that every later version remains unresolved \sref{30007658}{1}.

\par\textbf{Historical priority.} Earliest directly inspected qualifying passage in this audit; earlier origins were not established. A review or research-agenda statement does not imply original priority.

\subsection{Short English statement}
In the explicitly specified setting: Which combinations of enforcement, registration incentives, public services, and social insurance make formalization raise productivity and welfare rather than merely displace informal work? The mathematical statement above fixes the target, assumptions and scope of this catalogue formulation.

\subsection{Sources}
\begin{itemize}\raggedright
\item[\textnormal{[S1]}]\hypertarget{source:30007658:1}{} Gabriel Ulyssea, Matteo Bobba, Lucie Gadenne, Mariaflavia Harari. 2025. Informality, VoxDevLit 6(2). Chapter 6, Final Remarks, paragraphs 3–6.
\url{https://voxdev.org/voxdevlit/informality-issue-2/final-remarks-informality}.
{\small\textit{Earliest verified question/agenda reference; historical priority qualified below. Firm dynamics and workers’ formal-informal transitions are explicitly unresolved, including stepping-stone versus trap mechanisms.}}
\end{itemize}
\end{document}
